\section{Related Work}
\begin{remark} [Autonomous Agents for Web and Mobile Applications.]
Considerable initiatives have been invested in automating web navigation, driven by a vision of facilitating effortless human-web interaction.
Yet, previous research has been limited by the types of tasks and websites it can handle, either confined to simplified simulation environments~\cite{shiWorldBitsOpenDomain2017, Liu2018ReinforcementLO, yaoWebShopScalableRealWorld2022}, or limited to a narrow set of domains and websites~\cite{yaoWebShopScalableRealWorld2022, Xu2021GroundingOI}.
Recent studies~\cite{sunMETAGUIMultimodalConversational2022, Burns2022ADF, Li2020MappingNL} have utilized similar techniques for mobile applications, however, these are often simpler and offer fewer functions compared with full-fledged websites.
In contrast, \mind\ aims to adapt to a realistic web environment, characterized by its high diversity.
Also related is the research on web automation systems~\cite{leshedCoScripterAutomatingSharing2008, puppeteer}.
These technologies often demand programming skills, which can make them less accessible to general users. We aim to equip the web automation system with a natural language interface, thereby reducing the entry barrier significantly.
\end{remark}

\begin{remark} [Large Language Models.]
In recent years, there has been a surge in the development and application of large language models (LLMs). These models, often encompassing billions of parameters, are pre-trained on massive corpora of text data~\cite{Zhou2023ACS, LLMSurvey, Bommasani2021FoundationModels}, enabling them to capture intricate linguistic patterns, nuances, and relationships, resulting in unprecedented performance on a wide array of NLP tasks.
One of the most noteworthy attributes of LLMs is their few-shot learning capability. Unlike traditional machine learning models that necessitate extensive labeled data for task-specific fine-tuning, LLMs can often perform tasks with minimal task-specific examples. Furthermore, LLMs such as GPT-3~\cite{brownLanguageModelsAre2020} and PaLM~\cite{chowdhery2022palm} have also demonstrated the ability to do in-context learning, where they can adapt to novel tasks by simply providing context within the input prompt, eliminating the need for explicit retraining.
In this work, we explore the use of LLMs to build generalist agent on top of \mind\ by either tuning medium-sized LMs with only around \num{1000} examples, or prompting an LLM such as GPT-4, and have observed promising results.

\end{remark}

\begin{remark} [Grounded Language Understanding.]
Our work also aligns with the field of grounded language understanding, which aims to map natural language utterances onto executable plans in a target environment~\cite{pangu}. 
Many studies have centered around environments underpinned by a well-structured schema or ontology, including relational databases~\cite{spider, wangRATSQLRelationAwareSchema2020} and knowledge bases~\cite{yihSemanticParsingStaged2015,gu2021beyond}, which may not adequately reflect the more heterogeneous conditions in real-world situations.
Our work instead grounds natural language in the noisy and schemaless web environment.
Our setting is also connected to embodied AI, where an agent, guided by language instructions, carries out tasks in a physical environment~\cite{alfred, SayCan, llm-planner}.
Nonetheless, existing research primarily focuses on a specific setting (e.g., household environments), limiting their diversity.
\mind\ provides a unique testbed for studying a broad range of grounding challenges in real-world environments.
\end{remark}


\begin{remark} [Tool Learning.]
Recent developments have underscored LLMs' potential in using a myriad of tools (\ie, taking actions) to augment their capacity~\cite{augmented, tool-survey}, including search engine, translator, calculator, etc.
Example works include Toolformer~\cite{toolformer}, ReAct~\cite{react}, and ToolkenGPT~\cite{toolkengpt}.
The creation of recent benchmarks on tool learning~\cite{api-bank, gorilla} further highlights the growing interest in evaluating LLMs' proficiency in tool usage.
However, existing research primarily concentrates on short-term tool invocation, neglecting long-term planning. 
\mind\ can bridge this lacuna by necessitating LLMs to take actions within realistic web-browsing environments that demand prolonged decision-making sequences. 
Furthermore, \mind\ may stimulate the development of more advanced tools based on LLMs that interface the web with natural language\nop{, thereby obfuscating the low-level intricacies such as handling a DOM object}. 
These advanced tools could be subsequently employed by another LLM for more challenging problem-solving tasks~\cite{openagi, hugginggpt}.
\end{remark}